\section{Introduction}

\subsection{Motivation}

\subsection{Tasks and Requirements}

The scope of this work is categorized into four primary technical outputs: Robot Setup and Scenario Execution, Dataset Creation, Ground Truth Labeling, and Benchmark Criteria and Evaluation.

\subsubsection{Robot Setup and Scenario Execution}
This phase focuses on delivering fully operational platforms and execution scripts. It requires the successful software bring-up and deployment of the target mobile robots within the test environment. A precise characterization of each robot's embodiment is essential, capturing physical parameters (such as mass, geometry, footprint, and center of mass), kinematic limits (maximum velocity, acceleration, and braking distance), and sensor payload configurations. To ensure consistent replay across multiple runs, deterministic trajectory execution routines must be implemented. Safe operation is guaranteed through the integration of hardware and software override protocols, including e-stops and speed governors, to prevent collisions during dynamic testing. Finally, scenarios are parameterized and structured into families that systematically vary environmental and dynamic factors.

\subsubsection{Dataset Creation}
The primary deliverables for this stage are the raw and processed dynamic recordings. This entails acquiring a comprehensive multimodal dataset that spans all execution runs and parameter variations to ensure complete scenario coverage. Sensor streams—including LiDAR, RGB-D, IMU, and wheel odometry—must be standardized into well-documented ROS bag or MCAP formats, maintaining precise temporal synchronization and a structured data schema throughout the dataset.

\subsubsection{Ground Truth Labeling}
This output yields the annotated dataset records. It involves the development of a systematic ground-truth labeling pipeline to annotate high-precision spatial positioning, bounding box volumes, dynamic trajectories, and state classifications across all recorded datasets.

\subsubsection{Benchmark Criteria and Evaluation}
The final phase delivers the evaluation framework and associated metrics. This requires formulating standardized quantitative metrics for risk scoring to assess embodiment-dependent safety, maneuver success, and risk thresholds. These metrics will enable cross-embodiment comparative analysis across differing robotic platforms and environmental configurations.

\subsection{Contributions}

\subsection{Outline of the Report}
